\section{Comparison inputs, uncertainties and source limitations in detail}
\label{sup:detail:inputs}\label{sup:detail:numerics}\label{sup:detail:ensemble}
This section gives the inputs of each measured-curve comparison and the source statements behind them. The comparisons carry three kinds of uncertainty, which are not combined into one error bar: the reading half-width of a point digitized from a published trace; a source statement (mass basis, particle shape, probe depth, gas flow) that admits more than one reading, each an alternative model input, not a sample from a distribution; and the covariance or prediction interval of a coefficient fitted to independent data, which does not cover a temperature extrapolation or an unmeasured geometry. The main text states the strongest condition beside each headline result.

\subsection{Which source statements set a calculation}
For the 2008 thesis, the nominal sample mass is read as the wet charge, because the printed flux definition then reproduces the plotted constant-rate slope; reading it as dry meal predicts much less duty and a later transition, and the two readings cannot be averaged. Within a source, a tabulated quantity takes precedence over a position digitized from a graph, and a later peer-reviewed statement of a condition over an earlier thesis sentence. On version two of the digitization, the thesis sunflower window holds fourteen points from the measured crossing on. The earlier $9.1$-percent residual on the 2019 sunflower sphere belongs to version one ($14$ of $15$ in band); rescored on version two it is $11.6$ percent ($13$ of $14$), a change of the scored reading, not of the equations. For the 2019 run, the particle sizes and sphericity come from the article; the current experiment declares one representative radius because no mass histogram of the charge is printed. The film exposure multiplier comes from the constant-rate slope; the gas flow through the holder, the wetness near the thermocouple and its depth remain unmeasured (Sec.~\ref{sup:journalmarch}).

\subsection{External heat coefficient and finite sample geometry}
\label{sup:detail:geometry}
\paragraph{Film exposure multiplier: source data and calculation.}
\label{sup:fanermultiplier}
The multiplier was inferred from the constant-rate evaporation duty, including the Ackermann correction. For a
sphere with a wet surface,
$m''\Delta h_{\mathrm{vap}}=h_Q\Phi(\beta)(T_g-T_b)$ and
$r_c=3m''/(\rho_{\mathrm{dm}}\Rp)$, where $r_c$ is the decline of
hexane loading per unit dry mass and $\beta=m''c_{p,\mathrm v}/h_Q$.
The original construction matched mass-weighted particle rates to a
least-squares source slope of $7.26\times10^{-3}$ kg/(kg dry s), using
the four soybean solvent-loading readings in Fig.~2 of
\citet{faner2019kinetics}, above loading 0.2388. The digitized pairs
$(t_j\text{ [s]},X_j\text{ [kg/kg dry]})$ are
\[
\begin{aligned}
 &(0.000,\;0.492364),\quad(9.884,\;0.420900),\\
 &(20.077,\;0.337384),\quad(29.961,\;0.278064).
\end{aligned}
\]
These values keep the digitization precision for reconstruction; it does
not represent experimental accuracy.
With bars denoting the four-point arithmetic means, the slope is
\[
r_c=-\frac{\sum_j(t_j-\bar t)(X_j-\bar X)}
                  {\sum_j(t_j-\bar t)^2}
    =7.260603\times10^{-3}\ \mathrm{kg/(kg\ dry\ s)}.
\]
The construction used a truncated lognormal number distribution of sizes,
consistent with the reported mean diameter and size fractions and converted
to mass weights, with reference radius 0.666 mm and initial loading 0.4924.
For each size class $i$, the Ackermann balance gives
\[
r_i(f)=\frac{3f h^0_{Q,i}}{\rho_{\mathrm{dm}}R_i c_{p,\mathrm v}}
 \ln\!\left(1+\frac{c_{p,\mathrm v}(T_g-T_b)}
 {\Delta h_{\mathrm{vap}}}\right),
\]
where $h^0_{Q,i}$ is the unscaled correlation coefficient.
The earlier ensemble uses 400 mass-weighted classes. It chooses the
single factor $f$ so that the least-squares decline of its mean loading,
$\bar X(t;f)=\sum_i w_iX_i(t;f)$, at the same four times equals $r_c$.
The smallest classes enter the falling-rate period inside that window;
matching only the ensemble's initial rate would give 0.14508.
The four-time rule gives $f=0.1457854364$; the present 0.885-mm particle
uses it unchanged, with initial loading 0.475.
The multiplier is therefore a boundary calibration carried over from an
earlier geometry, not an independently measured transport property or
exposed surface fraction. The later falling-rate and temperature readings
were not used to identify it.
The digitized data, distribution inputs, slope calculation and exact
carried-over value are recorded in the reproducibility file
\texttt{faner\_film\_boundary\_derivation\_2026-10-08.json}.

The unscaled gas-film heat and mass coefficients use an industrial
reference superficial velocity of 0.230 m/s, obtained from a 3.9 kg/s gas
flow, a 6-m tower diameter, a gas density of 0.60 kg/m$^3$ and a bed voidage
of 0.40 \citep{coletto2022modeling}. At the present reference radius, the
heat coefficient is 104 W/(m$^2$ K) and the mass coefficient 0.0672 m/s;
scaling gives 15.1 W/(m$^2$ K) and 0.00980 m/s. Both coefficients
are scaled, while the geometric area remains $4\pi\Rp^2$.
The heat-duty inference does not determine mass transfer separately,
so applying the same factor to both is a declared boundary assumption.
Faner's charge of about 30 g in a 0.011-m$^2$ holder forms a layer
about three to five particle diameters deep \citep{faner2019kinetics}.
The industrial reference coefficients cannot be assumed to describe that
holder without a separate assessment of its gas flow and sample contact.

An unscaled calculation, with multiplier one and otherwise the same
four-cell particle, initial moisture, oil fraction, transport coefficients
and prescribed vapour, depletes external free hexane at 6.8 s and crosses
0.20 kg/kg dry at 6.8 s, rather than 46.8 s and 46.6 s, respectively.
Its heat and mass coefficients are 6.86 times larger. The calculation
enters regime B and continues to 6.88 s with the same component
and energy checks; the bounded continuation then reaches its computing-time
limit. This result shows the early boundary sensitivity; it is not a
complete falling-rate history or a prediction of the later temperature.

Separately, at the 2008 rig's own state, Whitaker's packed-bed correlation
predicts $0.82$--$1.06$ of the constant-rate duties on the wet-charge
basis (Sec.~\ref{sup:filmmatch}). Two inferred coefficients lie $3.7$ and
\SI{4.2}{\percent} below it, so the bracketing criterion over three
correlations, declared before the comparison, is narrowly missed; the film
is reported as a duty comparison with the departure stated. A tray
comparison needs its own gas composition, voidage, interfacial area and
condensation balance; the holder's coefficient cannot simply be copied.

\subsection{The declared spheres and the effective diffusivity law}
This subsection details the declared spheres and diffusivity law of Sec.~\ref{sup:surfaceloading}. The convention $R_d=\psi d_{\mathrm{eq}}/2$ gives the sphere the particle's volume-to-surface ratio. It was adopted after the effective-radius sensitivity had been inspected and is an assumption for the 2008 samples, whose shape was not measured; the 2019 sphericity $0.74$ comes from later samples. At constant diffusivity the spherical diffusion time scales with $\Rp^2/\Deff$, so a drying curve alone cannot separate radius from diffusivity. The volume-to-surface equivalence is an early-time construction, while the scored window reaches a Fourier number near $0.3$. The law's temperature term is weakly resolved (the 95-percent half-width of the activation energy exceeds its central value), and its prediction factor is a spread within the regression domain, not a statement about the \SI{41}{\kelvin} transfer above it. Reading the coefficient at a weighted mean concentration instead of at the surface would raise it five- to seventy-five-fold and dry far too fast; reading it in each interior cell gave no mesh-refined solution.

Table~\ref{tab:detail:geometry} compares these sphere conventions with
the same coefficient law. Errors are RMS fractions of the measured loading
drop; parentheses count points within the reading band, out of fourteen sunflower and fifteen soybean observations.

\begin{table}[htbp]\centering\footnotesize
\caption{Particle geometry sensitivity in the 2008 falling-rate comparisons.}
\label{tab:detail:geometry}
\begin{tabular}{@{}>{\raggedright\arraybackslash}p{0.49\textwidth}cc@{}}\toprule
Sphere convention & Sunflower & Soybean\\\midrule
Thesis volume-equivalent diameter & $0.379$ ($1$) & $0.310$ ($1$)\\
2019 soybean volume-equivalent diameter, both traces & $0.317$ ($1$) & $0.303$ ($1$)\\
$\psi$ times thesis sample diameter, declared & $0.172$ ($9$) & $0.156$ ($8$)\\
$\psi$ times 2019 soybean diameter, both traces & $0.099$ ($14$) & $0.147$ ($9$)\\
$\psi$ times 2019 sunflower diameter, sunflower only & $0.116$ ($13$) & --\\
\bottomrule\end{tabular}\end{table}

Table~\ref{tab:detail:dryladder} tests mesh and time-step sensitivity
on each thesis sample's declared sphere. It reports the RMS error of the
surface-read Fickian calculation as a fraction of the scored loading drop.

\begin{table}[htbp]\centering\footnotesize
\caption{Space and time refinement of the 2008 Fickian comparisons.}
\label{tab:detail:dryladder}\label{tab:evidence:dryladder}
\begin{tabular}{@{}lccc@{}}\toprule
Trace & 60 cells, nominal step & 120 cells, quarter step & 240 cells, quarter step\\\midrule
Sunflower, $\Rp=0.7215$ mm & $0.1721$ & $0.1678$ & $0.1676$\\
Soybean, $\Rp=0.6660$ mm & $0.1562$ & $0.1525$ & $0.1524$\\
\bottomrule\end{tabular}\end{table}

Across the law's prediction band, the residual on the declared sunflower sphere runs from $5.4$ to $52.6$ percent and on soybean from $7.7$ to $42.2$ percent; the best match lies inside the band, so points inside a broad band do not identify a mechanism. The central marches dry more slowly than both traces. Even so, a storage-consistent path with no coefficient fitted to the falling-rate points comes considerably closer than the same law on the volume-equivalent spheres. The uncertainty is dominated by particle shape and source conditions, not by truncation error in the surface-read form.
